\BourbakiAppendixExercises{Exercises}

\begin{enumerate}
\def\labelenumi{\arabic{enumi})}
\item
  Give an example of an invertible matrix A over a field D whose transpose is not invertible, and of an invertible matrix B over D whose transpose is invertible but satisfies \(\det^{t}B \neq \det B\) (consider square matrices of order 2 over a quaternion field).
\item
  Let D be a local ring; we denote by \(D^{*}\) the group of invertible elements of D, by \(D_{ab}^{*}\) the quotient of \(D^{*}\) by its derived group, and by \(\pi : D^{*} \to D_{ab}^{*}\) the canonical homomorphism. Let V be a free right D-module of finite dimension n.~We denote by \(\mathrm{B(V)}\) the set of bases of V and by \(\Omega(\mathrm{V})\) the set of mappings \(\omega : \mathrm{B(V)} \to \mathrm{D}_{\mathrm{ab}}^{*}\) satisfying the following two conditions:
\end{enumerate}

\begin{enumerate}
\def\labelenumi{(\roman{enumi})}
\item
  For \(\lambda_1, \ldots, \lambda_n\) in D* and \((v_1, \ldots, v_n) \in \mathrm{B}(\mathrm{V})\) , we have \(\omega(v_1 \lambda_1, \ldots, v_n \lambda_n) = \omega(v_1, \ldots, v_n)\pi(\lambda_1 \ldots \lambda_n)\) .
\item
  For \(i \neq j\) and \(\lambda \in D\) , we have \(\omega(v_{1},\ldots,v_{i}+v_{j}\lambda,\ldots,v_{n}) = \omega(v_{1},\ldots,v_{n})\) . a) Let e be an element of V and W a submodule of V supplementary to eD. Let \(\varphi \in \Omega(\mathrm{W})\) ; prove that there exists a unique element \(\omega\) of \(\Omega(\mathrm{V})\) satisfying
\end{enumerate}

\[
\omega (w _ {1}, \dots , w _ {n - 1}, e) = \varphi (w _ {1}, \dots , w _ {n - 1})
\]

for every basis \((w_{1},\ldots ,w_{n - 1})\) of W (adapt the proof of Proposition 1).

\begin{enumerate}
\def\labelenumi{\alph{enumi})}
\setcounter{enumi}{1}
\item
  Deduce that given two elements \(\omega\) and \(\omega'\) of \(\Omega(\mathrm{V})\) , there exists a unique element t of \(D_{ab}^{*}\) such that we have \(\omega' = t\omega\) .
\item
  Let u be an automorphism of V. Prove that there exists a unique element of \(D_{ab}^{*}\) , denoted by \(\det u\) , such that we have \(\omega(u(v_{1}),\ldots,u(v_{n}))=(\det u)\omega(v_{1},\ldots,v_{n})\) for every basis \((v_{1},\ldots,v_{n})\) of V and every element \(\omega\) of \(\Omega(\mathrm{V})\) . The mapping \(u\mapsto\det u\) from \(\mathbf{GL}(\mathrm{V})\) to \(D_{ab}^{*}\) is a group homomorphism.
\end{enumerate}

\(d\) ) Take \(\mathrm{V} = \mathrm{D}_d^n\) , so that \(\mathbf{GL}(\mathrm{V})\) is identified with the group \(\mathbf{GL}_n(\mathrm{D})\) of invertible square matrices of order \(n\) . Prove that the mapping \(\operatorname{det}:\mathbf{GL}_n(\mathrm{D})\to \mathrm{D}_{\mathrm{ab}}^*\) is the unique group homomorphism with value 1 on all matrices \(B_{ij}(\lambda)\) ( \(i\neq j\) , \(\lambda \in \mathrm{D}\) ) and satisfying \(\operatorname{det}(\operatorname{diag}(\lambda_1,\dots ,\lambda_n)) = \pi (\lambda_1\dots \lambda_n)\) for \(\lambda_1,\ldots ,\lambda_n\) in \(\mathrm{D}^*\) .

\begin{enumerate}
\def\labelenumi{\alph{enumi})}
\setcounter{enumi}{4}
\tightlist
\item
  Adapt Examples 2 through 6 of No.~4 to this setting.
\end{enumerate}

¶ 3) Let \(\mathrm{A} = \mathrm{A}_0 \oplus \mathrm{A}_1\) be a graded ring of type \(\mathbf{Z}/2\mathbf{Z}\) . Suppose that \(\mathrm{A}_0\) is contained in the center of \(\mathrm{A}\) and that every element of \(\mathrm{A}_1\) has square zero.

\begin{enumerate}
\def\labelenumi{\alph{enumi})}
\item
  Let \(A_{1}^{2}\) be the Z-submodule of \(A_{0}\) generated by the products of two elements of \(A_{1}\) ; show that \(A_{1}^{2}\) is a nil ideal of \(A_{0}\) and that \(A_{1}^{2} \oplus A_{1}\) is a two-sided nil ideal of A.
\item
  Let \(p, q\) be integers \(\geqslant 0\) . Denote by \(A^{p|q}\) the free graded right A-module \(A_d^p \oplus A_d(1)^q\) (II, §11, No.~2, p.~365, Example 3) and by \(\mathbf{GL}(p|q, \mathrm{A})\) the group of (graded) automorphisms (of degree 0) of \(A^{p|q}\) . Such an automorphism is represented in the canonical basis of \(A^{p|q}\) by a matrix \(X = \begin{pmatrix} A & B \\ C & D \end{pmatrix}\) with \(A \in \mathbf{M}_p(\mathrm{A}_0)\) , \(B \in \mathbf{M}_{p,q}(\mathrm{A}_1)\) , \(C \in \mathbf{M}_{q,p}(\mathrm{A}_1)\) , \(D \in \mathbf{M}_q(\mathrm{A}_0)\) . Prove that the matrices \(A\) and \(D\) are invertible (use \(a\) ) and Exercise 16 of VIII, p.~168).
\item
  For X as above, set \(\operatorname{sdet} X = \det(A - BD^{-1}C) \det D^{-1}\) . Prove that \(\operatorname{sdet} X\) is an invertible element of \(A_{0}\) .
\end{enumerate}

\(d\) ) Prove that every matrix \(X\in \mathbf{GL}(p|q,\mathrm{A})\) can be written as a product

\[
X = \left( \begin{array}{c c} I & R \\ 0 & I \end{array} \right) \left( \begin{array}{c c} P & 0 \\ 0 & Q \end{array} \right) \left( \begin{array}{c c} I & 0 \\ S & I \end{array} \right).
\]

(Take \(R = BD^{-1}\) , \(P = A - BD^{-1}C\) , \(Q = D\) , \(S = D^{-1}C\) .)

\begin{enumerate}
\def\labelenumi{\alph{enumi})}
\setcounter{enumi}{4}
\tightlist
\item
  Prove that sdet is a homomorphism from \(\mathbf{GL}(p|q,\mathrm{A})\) to the group of invertible elements of \(A_{0}\) . (Use d) to reduce to proving the equality \(\operatorname{sdet}(XY)=\operatorname{sdet}X\operatorname{sdet}Y\) when X is of the form \(\binom{I}{S}\operatorname{sdet}(I)\) , where S is an elementary matrix. Observe that for every matrix \(R\in\mathbf{M}_{p,q}(\mathrm{A}_{1})\) , the product of two arbitrary entries of the matrix RS (resp. SR) is zero, and deduce the equality \(\det(I-RS)=(\det(I-SR))^{-1}\) .
\item
  Let \(X = \begin{pmatrix} A & B \\ C & D \end{pmatrix}\) be an element of \(\mathbf{GL}(p|q,\mathrm{A})\) ; set \(^{st}X = \begin{pmatrix} t_{A} & t_{C} \\ -t_{B} & t_{D} \end{pmatrix}\) and \(^{\pi}X = \begin{pmatrix} D & C \\ B & A \end{pmatrix}\) . Let \(Y \in \mathbf{GL}(p|q, A)\) ; we have \(^{st}(XY) = ^{st}Y^{st}X\) and \(^{\pi}(XY) = ^{\pi}X^{\pi}Y\) . Prove the equalities \(\operatorname{sdet}X = \operatorname{sdet}^{st}X = (\operatorname{sdet}^{\pi}X)^{-1}\) (use d) to calculate \(\operatorname{sdet}^{\pi}X\) .
\end{enumerate}

\begin{enumerate}
\def\labelenumi{\arabic{enumi})}
\setcounter{enumi}{3}
\tightlist
\item
  Keep the assumptions of Exercise 3. Let \(\mathrm{M} = \mathrm{M}_0 \oplus \mathrm{M}_1\) be a finitely generated free graded right A-module; there exist unique integers \(p, q\) such that \(\mathrm{M}\) is isomorphic to \(\mathrm{A}^{p|q}\) . We say that a basis of \(\mathrm{M}\) is adapted if its first \(p\) prime vectors are homogeneous of degree 0 and the last \(q\) are homogeneous of degree 1.
\end{enumerate}

\begin{enumerate}
\def\labelenumi{\alph{enumi})}
\item
  Let u be an automorphism of the graded module M, and let \(M_{\mathcal{B}}(u)\) be its matrix with respect to an adapted basis of B. Prove that the element sdet \(M_{\mathcal{B}}(u)\) does not depend on the choice of B. We denote it by sdet u.
\item
  Denote by \(u(1)\) the endomorphism u viewed as an automorphism of the graded module M(1). If B is an adapted basis of M, then the basis \(^{\pi}B\) obtained by interchanging the first p vectors of B and the last q is an adapted basis of M(1), and we have \(M_{\pi\mathcal{B}}(u(1)) = ^{\pi}M_{\mathcal{B}}(u)\) (Exercise 3, f)). Deduce that we have \(\operatorname{sdet}(u(1)) = (\operatorname{sdet} u)^{-1}\) .
\item
  On every graded left A-module N, we define a canonical graded right A-module structure by setting \(ax = (-1)^{ij}xa\) for \(a \in A_i\) , \(x \in N_j\) . In particular, we view the graded dual \(M^{*gr}\) of M as a graded right A-module, also isomorphic to \(A^{p|q}\) . Let u be an automorphism of M. We denote by \({}^{st}u\) the automorphism of \(M^{*gr}\) characterized by the formula \(\langle{}^{st}u(x^{*}), x\rangle = (-1)^{rs}\langle x^{*}, u(x)\rangle\) for \(x^{*} \in (\mathrm{M}^{*\mathrm{gr}})_r\) , \(x \in M_s\) . If B is an adapted basis of M, then the dual basis \(B^*\) is an adapted basis of \(M^{*gr}\) , and we have \(M_{\mathcal{B}^*}({}^{st}u) = {}^{st}M_{\mathcal{B}}(u)\) (Exercise 3, f). We have \(\operatorname{sdet}({}^{st}u) = \operatorname{sdet}(u)\) .
\end{enumerate}

